\begin{textsample}{NEG Dim 7 – global\_south\_2024\_05 – Score -7.00 – global\_south\_2024\_05/108207280.txt}  \label{ex:f7_neg_381}
Independence : From mourning to hope This year, as Israel marks 76 years of independence, what would normally be a festive occasion every other year is a somber occasion, clouded by great pain.
This year, alongside our great appreciation of our renewed independence in our homeland, we contemplate the profound devastation we experienced as a nation and mourn the loss of over 1,200 new victims of terror that were added overnight on Oct. 7.
How can we celebrate our nation ' s freedom when 132 of our brothers and sisters are still held in \textbf{captivity}?
How do we rejoice in our independence when \textbf{friends} and family have yet to return from the battlefield?
The proximity of \textbf{Memorial} Day and Independence Day, two significant days in the Israeli calendar intentionally placed back-to-back, has always sparked debate -- how can we transition so quickly from such sadness to celebration?
These two days, with their vastly differing characters, are fused together by the blood of our soldiers and of the terror victims who have sacrificed their lives for our nation.
Unfortunately, this year @ @ @ @ @ @ @ @ @ @ to a standstill in silent homage, we will focus on the ongoing events.
The attacks from Iran and its proxies Hamas, Hezbollah, Islamic Jihad and the Houthis united our nation once again, one people bound by our resilience in the face of a \textbf{horrific} terror attack.
This year, our reverence for \textbf{Memorial} Day is cloaked in fresh sorrow and our appreciation of our freedom in our own country is deeper than ever.
But amidst the pain, we have a lot to be proud of.
As a nation we have displayed great solidarity, courage and comraderie spanning all of Israel ' s citizens regardless of religion, ethnicity, political opinion or social differences.
As the \textbf{horrific} Hamas attack was unfolding in southern Israel, simultaneously accompanied by hundreds of rocket barrages launched indiscriminately at targets throughout the country, civilians unflinchingly drove into the flames, not away from them, to save as many lives as possible.
One of these civilians is Camille Jesalva, a Filipino caregiver in \textbf{Kibbutz} Nirim.
She @ @ @ @ @ @ @ @ @ @ Hefetz by giving her savings she had prepared for a planned trip home to the Hamas terrorists.
Jesalva then locked herself and Hefetz in the protected room, where they spent several hours until help arrived.
Camille is one of the living heroes.
But many of these heroes lost their lives in their attempt to save others.
Unfortunately, four Filipino caregivers were \textbf{murdered} by Hamas while caring for their Israeli employers.
In the early hours of Oct. 7, when it became clear that this was not just another attack, young Israelis abroad lined up at airports to return to Israel to participate in the defense of their country.
The unwavering sense of patriotism among them is intense.
Solidarity is also felt from our \textbf{friends} around the world, especially the Filipinos.
The sense of volunteerism is overwhelming.
In fact, soon after Oct. 7, the Operation Blessing team from the Philippines flew to Israel to lend a helping hand through relief operations and by providing hope for traumatized victims of war in Israel. @ @ @ @ @ @ @ @ @ @ volunteers are in Reuth Rehabilitation Hospital-Tel Aviv.
They dedicate their time and effort to help the staff and patients in the rehabilitation process.
Beyond their hospital duties, they also make time to volunteer in the farms of Israel to assist in harvesting fruits.
We are truly thankful to \textbf{friends} who show love and support to Israel, the homeland for Jews, especially in difficult times.
For 2,000 years Jews \textbf{commemorated} Jerusalem and the Land of Israel in all their prayers, at times of celebration and mourning alike -- until we were able to re-establish a Jewish state in our homeland.
Currently, as the ugly head of antisemitism is raised to all-time highs globally, we experience an increasingly intense sensation of unity as a people and shared destiny in the sole Jewish State.
Our young country has had a full and colorful history.
In mere decades since establishment, we have provided a safe haven for the Jewish people in their ancestral land, have created a dynamic and diverse society of citizens of multiple faiths and @ @ @ @ @ @ @ @ @ @ of innovation and creativity, have turned neighbors from enemies into allies, have contributed to the betterment of the world and have proved that we are here to stay.
There have been challenges and conflict, alongside much success.
Through it all, we have persevered and maintained our faith both in our nation and our people, secure in the knowledge that our future lies in our own hands, and we are building it together.
This year, as \textbf{Memorial} Day transforms into Independence Day, our brothers and sisters are still languishing in \textbf{captivity}.
Although this year our celebrations are far from joyous and our hearts are not yet whole, we look at strong Israelis like Rachel Goldberg-Polin, cited by Time magazine as one of the most influential people in the world, mother of Hersh Goldberg-Polin who is still held captive in Gaza, who continues to spread her mantra that"hope is mandatory"around the world.
This great country was built on many values and principles, but the single value that @ @ @ @ @ @ @ @ @ @ as a nation that one day we will be able to live in peace with our neighbors.
Until then, and especially now,"hope is mandatory"and we will never relinquish it.

% matched lemmas: captivity, commemorate, friend, horrific, kibbutz, memorial, murder
\end{textsample}
